% !TEX program = xelatex
% متن فارسی و قلم همراه؛ موتور کامپایل باید XeLaTeX باشد.
\documentclass[11pt,a4paper]{article}
\usepackage[margin=2.2cm,headheight=16pt]{geometry}
\usepackage{graphicx,booktabs,amsmath,array,float,pdflscape}
\usepackage{xcolor}
\definecolor{accent}{HTML}{187B8D}
\usepackage[colorlinks=true,urlcolor=accent,linkcolor=accent]{hyperref}
\usepackage{xepersian}
\settextfont[
    Path=fonts/,
    BoldFont=Vazirmatn-Bold.ttf
]{Vazirmatn-Regular.ttf}
\setlatintextfont{Latin Modern Roman}
\graphicspath{{../../figures/}{figures/}}
\newcommand{\TrainCount}{5382}
\newcommand{\ValCount}{1232}
\newcommand{\TestCount}{1359}
\newcommand{\LowTrain}{434}
\newcommand{\ClipTrain}{101}
\newcommand{\MajorityPct}{12.76}
\linespread{1.4}
\setlength{\parskip}{5pt}
\setlength{\parindent}{0pt}
\newcommand{\fig}[3]{\begin{figure}[H]\centering\includegraphics[width=#2\linewidth]{#1}\caption{#3}\end{figure}}
\begin{document}
\begin{titlepage}
\vspace*{1.1cm}
{\color{accent}\Large پروژهٔ درس تحلیل داده}\par
\vspace{1.2cm}
{\Huge\bfseries تشخیص فرمان‌های گفتاری}\par
\vspace{0.5cm}
{\Large برای گویندگان جدید}\par
\vspace{0.8cm}
{\large فاز اول: فهم مسئله، ممیزی کیفیت و تحلیل اکتشافی}\par
\vspace{1cm}
\rule{\linewidth}{1pt}
\textbf{نام دانشجو:} مهرشاد حقیقت\\
\textbf{شماره دانشجویی:} ۴۰۲۱۰۰۴۱۸\\
\textbf{مسیر:} د؛ تحلیل صوت و تشخیص کلیدواژه\\
\textbf{داده:} \lr{Mini Speech Commands}، زیرمجموعهٔ رسمی آموزش \lr{TensorFlow}\\
\end{titlepage}

\section{خلاصه و تعریف مسئله}
در این پروژه بررسی می‌کنم که تشخیص فرمان کوتاه برای گوینده‌ای که در آموزش حضور نداشته تا چه اندازه قابل اعتماد خواهد بود. این مسئله می‌تواند به انتخاب روش مناسب برای یک رابط فرمان محدود کمک کند؛ ارزیابی کاربرد واقعی مستلزم بررسی گفتار ناشناخته، سکوت و شرایط ضبط متنوع‌تر است.

واحد تحلیل و پیش‌بینی یک کلیپ صوتی است. برچسب هدف یکی از هشت واژهٔ:
\begin{center}
    \lr{down, go, left, no, right, stop, up, yes}
\end{center} است. شناسهٔ گوینده صرفاً برای تقسیم و تحلیل تعمیم استفاده می‌شود و ورودی مدل نخواهد بود.
\begin{center}
\begin{tabular}{rr}
\toprule
شاخص & مقدار \\
\midrule
تعداد فایل خام & 8000 \\
موج یکتا برای تحلیل & 7973 \\
شناسهٔ گوینده & 1750 \\
کلاس & 8 \\
نرخ نمونه‌برداری & ۱۶۰۰۰ هرتز \\
قالب & تک‌کاناله، ۱۶ بیت \\
کلیپ کوتاه‌تر از یک ثانیه & 822 \\
کمینه و بیشینهٔ طول & ۰٫۴۲۷ و ۱٫۰۰۰ ثانیه \\
\bottomrule
\end{tabular}
\end{center}

هر کلاس در دادهٔ خام ۱۰۰۰ فایل دارد. با مقایسهٔ موج رمزگشایی‌شده، ۲۷ نسخهٔ اضافی شناسایی شد؛ ۷۹۷۳ موج یکتا برای تحلیل باقی می‌ماند. این توازن اولیه به معنی نمایندگی توزیع فرمان‌ها در کاربرد واقعی نیست.
\fig{01_class_balance.png}{0.9}{ممیزی ساختار کل زیرمجموعه، پیش از کنارگذاشتن تکراری‌ها.}

\newpage
\section{منشأ داده و حدود دامنه}
داده از آرشیو آموزشی رسمی \lr{TensorFlow} دریافت شده است \cite{tutorial}. این آرشیو گزیده‌ای از مجموعهٔ \lr{Speech Commands} است؛ مقالهٔ اصلی، هدف جمع‌آوری گفتار کوتاه برای شناسایی واژگان محدود و روش ارزیابی آن را توضیح می‌دهد \cite{warden}. متن \lr{README} داخل آرشیو در بسته حفظ شده است.

انتخاب نسخهٔ کوچک به دلیل هزینهٔ کمتر دانلود، پردازش و بازاجرای آزمایش‌ها انجام شد؛ انتخاب قبل از آموزش مدل بوده و بر مبنای بهبود معیار مدل تغییر نکرده است. این نسخه با مجموعهٔ کامل و نسخهٔ دوازده‌کلاسهٔ \lr{TFDS} یکسان نیست \cite{tfds}.

این داده شامل گفتار گردآوری‌شده است، نه آزمایش تصادفی دربارهٔ تأثیر نویز، لهجه یا ویژگی‌های افراد. تاریخ ضبط هر فایل، جنسیت، لهجه، دستگاه ضبط و هویت واقعی افراد در این نسخه برای تحلیل حاضر در دسترس نیست.

\textbf{ساختار:} آرشیو شامل پوشه‌های فرمان و فایل‌های \lr{WAV} است. نام فایل از الگوی شناسهٔ گوینده، عبارت \lr{\_nohash\_} و شمارهٔ تکرار پیروی می‌کند. استخراج شناسه، بررسی قالب نام و ساخت جدول فایل‌ها در کد انجام شده است. فایل‌های فرادادهٔ سیستم‌عامل وارد شمارش صوت نمی‌شوند.

\textbf{بازاجرا:} کد هنگام دریافت و هر بار اجرا، هش \lr{SHA-256} آرشیو را با مقدار ثبت‌شده در پیکربندی تطبیق می‌دهد. دادهٔ خام تغییر نمی‌کند. شاخص‌ها در جدول موجودی ذخیره می‌شوند و انتخاب تکراری‌ها و تقسیم‌ها از روی همین جدول قابل بررسی‌اند.

\subsection{سه سؤال پژوهشی و انتخاب سؤال اصلی}
\begin{enumerate}
\item طول، انرژی و شاخص‌های کیفیت چگونه بین فرمان‌ها و گویندگان تغییر می‌کنند؟ پاسخ اولیه در همین فاز بررسی می‌شود.
\item روش سادهٔ مبتنی بر ویژگی صوتی و مدل کوچک طیفی تا چه اندازه به گویندگان جدید تعمیم می‌یابند؟ این سؤال اصلی پروژه است.
\item نویز کنترل‌شده چه افتی در عملکرد ایجاد می‌کند و آیا افزایش داده آن را کاهش می‌دهد؟ این سؤال توسعه‌یافتهٔ فازهای بعد است.
\end{enumerate}
سؤال دوم انتخاب شد چون هدف پیش‌بینی و واحد ارزیابی مشخص دارد و مستقیماً به مقایسهٔ مدل، تحلیل خطا و استحکام متصل می‌شود. دادهٔ فاز اول هنوز پاسخ عملکردی به آن نمی‌دهد.

\newpage
\section{روش اندازه‌گیری و ممیزی کیفیت}
فایل‌های \lr{PCM} تک‌کانالهٔ ۱۶ بیتی به مقیاس ثابت تبدیل می‌شوند:
\begin{equation}
x_n=\frac{a_n}{32768},\qquad
\mathrm{RMS}=\sqrt{\frac{1}{N}\sum_{n=1}^{N}x_n^2},\qquad
L=20\log_{10}\bigl(\max(\mathrm{RMS},10^{-12})\bigr).
\end{equation}
مقیاس \lr{dBFS} نسبت به دامنهٔ کامل دیجیتال است و بلندی فیزیکی صدا بر حسب فشار صوت را اندازه نمی‌گیرد. قله، میانگین سیگنال، طول، تعداد نمونه و نسبت نمونه‌های صفر نیز ثبت شده‌اند.

\textbf{اشباع:} نمونه‌هایی برابر با دو حد \lr{PCM} یعنی منفی ۳۲۷۶۸ یا مثبت ۳۲۷۶۷ شمارش می‌شوند. نسبت حداقل ۰٫۱ درصد یک پرچم است؛ رسیدن به حد دیجیتال به‌تنهایی اثبات نمی‌کند فایل برای طبقه‌بندی نامناسب است.

\textbf{کم‌انرژی:} انرژی کل کمتر از منفی ۴۰ دسی‌بل پرچم می‌گیرد. قاب‌های کامل ۲۰ میلی‌ثانیه‌ای با انرژی کمتر از منفی ۴۵ دسی‌بل برای برآورد نسبت قاب کم‌انرژی استفاده می‌شوند؛ قاب انتهایی ناقص وارد این نسبت نمی‌شود. این شاخص جایگزین تشخیص گفتار، سکوت یا اندازه‌گیری نسبت سیگنال به نویز نیست.

\textbf{تکرار:} هش نمونه‌های رمزگشایی‌شده همراه نرخ نمونه‌برداری مقایسه می‌شود. از هر موج یکسان، اولین مسیر در ترتیب واژگانی برای تحلیل انتخاب می‌شود؛ سایر نسخه‌ها با پرچم در موجودی باقی می‌مانند. شباهت تقریبی، بازنمونه‌برداری یا تغییر دامنه ممکن است با این روش تشخیص داده نشود.
\begin{center}
\begin{tabular}{rrr}
\toprule
بررسی روی کل داده & تعداد & تصمیم اولیه \\
\midrule
خطای خواندن & 0 & هیچ مورد \\
نسخهٔ اضافی PCM & 27 & یک نسخه در تحلیل \\
تعارض برچسب تکراری & 0 & هیچ مورد \\
تمام نمونه‌ها صفر & 0 & هیچ مورد \\
کمتر از ۰٫۲۵ ثانیه & 0 & هیچ مورد \\
انرژی کمتر از منفی ۴۰ & 691 & ثبت و بررسی؛ حفظ \\
اشباع حداقل ۰٫۱ درصد & 152 & ثبت و بررسی؛ حفظ \\
\bottomrule
\end{tabular}
\end{center}

تمام ۸۰۰۰ فایل قابل رمزگشایی‌اند و همگی ۱۶ کیلوهرتز، تک‌کاناله و ۱۶ بیتی هستند. شناسه و برچسب استخراج‌شدهٔ خالی نداریم؛ این موضوع به معنی کامل‌بودن فرادادهٔ جمعیت‌شناختی نیست. ۵۴ فایل در ۲۷ گروه تکراری قرار دارند و تعارض برچسب در این گروه‌ها مشاهده نشد.

\newpage
\section{توزیع‌ها و تفاوت فرمان‌ها}
از این بخش به بعد، تحلیل اکتشافی بر \TrainCount{} فایل یکتای آموزش انجام می‌شود. فایل‌های اعتبارسنجی و آزمون فقط برای ممیزی ساختاری و کنترل مرز تقسیم بررسی شده‌اند؛ هیچ انتخاب مدل بر اساس آن‌ها انجام نشده است.
\fig{02_distributions.png}{1}{توزیع طول و انرژی در آموزش؛ تجمع طول‌ها در یک ثانیه و دنبالهٔ کم‌انرژی دیده می‌شود.}
بیشتر کلیپ‌ها یک ثانیه‌اند؛ ۵۱۰ کلیپ آموزشی کوتاه‌تر از یک ثانیه‌اند. طول کوتاه‌تر لزوماً خرابی نیست و می‌تواند با مدت تلفظ و برش فایل مرتبط باشد. برای فاز دوم، افزودن صفر به انتهای کلیپ تا طول یک ثانیه در برنامهٔ ادامهٔ پروژه قرار دارد؛ سیگنال خام در این فاز تغییر نکرده است.
\fig{03_class_quality.png}{1}{مقایسهٔ طول و انرژی بر حسب فرمان؛ نقاط بیرون سبیل‌ها به‌طور خودکار حذف نشده‌اند.}
میانهٔ انرژی آموزش حدود منفی ۲۶٫۷۵ دسی‌بل است. دامنهٔ بین‌چارکی انرژی و نقاط کم‌انرژی در تمام کلاس‌ها دیده می‌شود. شباهت توزیع انرژی میان فرمان‌ها نشان می‌دهد انرژی به‌تنهایی نشانهٔ تفکیک قطعی کلاس نیست؛ این برداشت جای ارزیابی مدل را نمی‌گیرد.

\newpage
\section{روابط چندمتغیره و حساسیت آستانه}
\fig{06_correlations.png}{0.65}{همبستگی رتبه‌ای اسپیرمن بین شاخص‌های صوتی در آموزش.}
همبستگی انرژی و قلهٔ دامنه حدود ۰٫۹۴ است؛ هر دو از بزرگی دامنه تأثیر می‌گیرند. همبستگی انرژی و نسبت قاب کم‌انرژی حدود منفی ۰٫۵۲ است. بخشی از این رابطه از تعریف مشترک مبتنی بر انرژی ناشی می‌شود؛ رابطهٔ علّی دربارهٔ کیفیت ضبط یا گوینده استنباط نمی‌شود. همبستگی طول و انرژی حدود ۰٫۰۵ است و رابطهٔ یکنواخت قوی نشان نمی‌دهد.
\fig{09_low_energy.png}{0.94}{کم‌انرژی‌بودن قاب‌ها و حساسیت تعداد پرچم‌های کلیپ به آستانه؛ فقط آموزش.}
با آستانه‌های منفی ۵۰، ۴۵، ۴۰ و ۳۵ دسی‌بل، به‌ترتیب ۱۰۵، ۲۰۸، ۴۳۴ و ۹۶۶ کلیپ آموزش علامت می‌خورند. وابستگی زیاد شمارش به آستانه دلیل دیگری برای پرهیز از حذف خودکار است. میانهٔ نسبت قاب‌های کم‌انرژی ۰٫۵۸ است؛ این عدد سهم سکوتِ تأییدشده نیست.

\newpage
\section{ممیزی گویندگان}
در کل داده ۱۷۵۰ شناسهٔ متمایز وجود دارد. در آموزش ۱۱۹۳ شناسه حضور دارند؛ میانهٔ تعداد کلیپ یکتا به ازای هر شناسه ۳ و بیشینه ۲۵ است. این شناسه‌ها نمایندهٔ گروه‌بندی فراهم‌شده‌اند؛ استقلال هویت واقعی انسان‌ها فراتر از شناسه‌های داده بررسی نشده است.
\fig{04_speakers.png}{1}{مشارکت نامتوازن گویندگان و تعداد فرمان‌های پوشش‌داده‌شده در آموزش.}
\begin{center}
\begin{tabular}{rrrrrr}
\toprule
فرمان & تعداد & گوینده & میانه انرژی & کم‌انرژی & اشباع \\
\midrule
\lr{down} & 670 & 466 & \lr{-26.27} & 55 & 18 \\
\lr{go} & 673 & 473 & \lr{-26.34} & 41 & 9 \\
\lr{left} & 660 & 474 & \lr{-26.72} & 53 & 12 \\
\lr{no} & 676 & 486 & \lr{-26.81} & 56 & 14 \\
\lr{right} & 687 & 499 & \lr{-26.54} & 61 & 10 \\
\lr{stop} & 672 & 466 & \lr{-26.78} & 54 & 14 \\
\lr{up} & 678 & 495 & \lr{-27.68} & 63 & 11 \\
\lr{yes} & 666 & 466 & \lr{-26.85} & 51 & 13 \\
\bottomrule
\end{tabular}
\end{center}

تعداد گویندهٔ هر سطر هم‌پوشان است و جمع این ستون تعداد کل گویندگان نیست. متوازن‌بودن تعداد فایل‌های خام به معنی توازن تعداد گویندگان یا کیفیت صوت نیست.

\newpage
\section{نمونه‌های موج و نمایش طیفی}
برای هر فرمان سه نمونه از آموزش با مرتب‌سازی هش مسیر و قاعدهٔ ثابت انتخاب شده است؛ نمونه‌ها برای زیباترشدن شکل یا جدایی بهتر کلاس‌ها دست‌چین نشده‌اند. فهرست دقیق در \lr{plot\_examples.csv} است. نمایش طیفی با پنجرهٔ \lr{Hann} با طول ۴۰۰ نمونه، هم‌پوشانی ۲۴۰ و \lr{FFT} با طول ۵۱۲ ساخته شده است. نرخ ۱۶ کیلوهرتز، بسامد تا ۸ کیلوهرتز را پوشش می‌دهد.
\fig{08_examples_1.png}{1}{سه نمونه از هر یک از چهار فرمان اول؛ موج و طیف توان با مقیاس رنگ مشترک.}
در نمونه‌های رسم‌شده، محل آغاز و پایان بخش پرانرژی یکسان نیست. برای مثال، نمونه‌های \lr{go} در جایگاه‌های زمانی متفاوتی فعال می‌شوند؛ در \lr{left} نیز دامنهٔ نمونهٔ سوم نسبت به نمونهٔ دوم کوچک‌تر است. یکسان‌نبودن دامنه و جایگاه زمانی، لزوم بررسی روش نرمال‌سازی و انتقال زمانی در فازهای بعد را نشان می‌دهد؛ به‌تنهایی شاهد بهبود مدل با آن روش‌ها نیست.

\newpage
\fig{08_examples_2.png}{1}{سه نمونه از هر یک از چهار فرمان بعدی؛ انتخاب از آموزش و بدون شنود انسانی ثبت‌شده.}
در \lr{stop} یک بخش پرانرژی نسبتاً متمرکز و در برخی نمونه‌های \lr{yes} محتوای پربسامد گسترده‌تر دیده می‌شود. این تفاوت‌های بصری فقط دربارهٔ نمونه‌های انتخاب‌شده‌اند؛ تعمیم آن‌ها به تمام تلفظ‌ها نیازمند بررسی بیشتر است. طیف کم‌رنگ می‌تواند ناشی از دامنهٔ کم باشد و تشخیص قطعی نبود گفتار نیست.
\fig{07_duration_energy.png}{0.8}{هم‌پوشانی فرمان‌ها در صفحهٔ طول و انرژی؛ اطلاعات زمانی و طیفی غنی‌تر برای مدل‌سازی لازم است.}

\newpage
\section{طراحی تقسیم و جلوگیری از نشت}
گروه پایه، شناسهٔ گوینده است. اگر موج یکسان بین دو شناسه باشد، هر دو به یک گروه متصل می‌شوند؛ در این نسخه چنین موردی مشاهده نشد. هش گروه همراه \lr{salt} ثابت، یک عدد در بازهٔ صفر تا یک می‌سازد. زیر ۰٫۷ به آموزش، از ۰٫۷ تا ۰٫۸۵ به اعتبارسنجی و بقیه به آزمون می‌رود. نتیجهٔ عددی با نسبت هدف یکسان نیست و برای رسیدن به نسبت یا معیار بهتر، \lr{salt} جست‌وجو نشده است.
\begin{center}
\begin{tabular}{rrrr}
\toprule
بخش & فایل خام & فایل تحلیل & گوینده \\
\midrule
آموزش & 5409 & 5382 & 1193 \\
اعتبارسنجی & 1232 & 1232 & 266 \\
آزمون & 1359 & 1359 & 291 \\
\bottomrule
\end{tabular}
\end{center}

\fig{10_splits.png}{0.95}{تعداد فایل خام هر کلاس در تقسیم‌ها؛ ۲۷ نسخهٔ اضافی در آموزش برای تحلیل کنار گذاشته شده‌اند.}
آزمون برنامه‌ای تأیید می‌کند هیچ شناسهٔ گوینده و هیچ هش موج بین تقسیم‌ها مشترک نیست. این تقسیم اختصاصی پروژه است و نباید با تقسیم رسمی بنچمارک مقایسهٔ مستقیم شود.
\begin{itemize}
\item مسیر فایل و شناسهٔ گوینده ورودی مدل نیستند؛ نام پوشه برچسب را آشکار می‌کند.
\item مقیاس‌بندی، \lr{PCA} و انتخاب ویژگی فقط روی آموزش برازش می‌شوند.
\item نسخه‌های افزوده‌شده از فایل، فقط در همان تقسیم فایل اصلی می‌مانند.
\item هنگام افزودن نویز واقعی، ضبط منبع پیش از برش به آموزش و ارزیابی جدا می‌شود.
\end{itemize}
آمار ساختاری آزمون در فاز اول دیده شده است؛ بنابراین ادعای دیده‌نشدن مطلق آزمون از آغاز مطرح نیست. از این مرحله، آزمون برای انتخاب مدل و آستانه استفاده نمی‌شود.

\newpage
\section{خط مبنا و برنامهٔ فازهای بعد}
خط مبنای ساده، پیش‌بینی کلاس اکثریت آموزش یعنی \lr{right} برای همهٔ ورودی‌هاست. سهم آن در آموزش ۷۶.۱۲ درصد است؛ این عدد نتیجهٔ آزمون یا برآورد قطعی عملکرد آینده نیست. در فاز سوم، این خط مبنا کنار رگرسیون لجستیک روی خلاصهٔ \lr{MFCC}، یک مدل غیرخطی و \lr{CNN} کوچک قرار می‌گیرد.

معیار اصلی در نظر گرفته‌شده \lr{macro-F1} است؛ دقت کلی، بازیابی هر فرمان و ماتریس درهم‌ریختگی نیز گزارش خواهند شد. برای عدم قطعیت، بازنمونه‌گیری در سطح گوینده از فرض استقلال تمام کلیپ‌ها مناسب‌تر است. انتخاب نهایی مدل و تنظیم پارامترها با اعتبارسنجی انجام می‌شود.

در فاز دوم، طول استاندارد و بازنمایی طیفی ایجاد و \lr{PCA} و خوشه‌بندی روی آموزش بررسی می‌شود. در فاز چهارم، خطاهای فرمان‌های مشابه، عملکرد گویندگان جدید و افت در چند شدت نویز تحلیل می‌شود. در فاز پنجم، حذف یک مؤلفهٔ ویژگی یا افزایش داده و سناریوهای تغییر شرایط به گزارش نهایی متصل می‌شوند. این موارد برنامه‌اند و در این فاز نتیجهٔ اجرایی برایشان ارائه نشده است.

\section{نویز، سکوت و محدودیت کاربرد}
آرشیو \lr{Mini} فایل مستقل نویز پس‌زمینه و کلاس \lr{unknown/silence} ندارد. با این حال، کلیپ‌های کم‌انرژی و قاب‌های کم‌انرژی ممیزی شده‌اند. نبود نمونهٔ تمام‌صفر اثبات نمی‌کند تمام کلیپ‌ها دارای فرمان قابل شنیدن‌اند. صف بازبینی برای شنیدن نمونه‌های نماینده و موارد مشکوک آماده است و هنوز وضعیت آن خالی است.

این پروژه در دامنهٔ فعلی طبقه‌بندی بستهٔ هشت فرمان است. در نتیجه، خطای مثبت یک فرمان در برابر دیگر فرمان‌ها قابل تعریف است، اما نرخ فعال‌شدن کاذب در سکوت یا گفتار ناشناخته با این مجموعه قابل تخمین نیست. برای ارزیابی واقع‌گرایانه‌تر باید منبع مجاز و مستقل نویز/گفتار خارج از دامنه اضافه و تغییر دامنه ثبت شود.

\section{اخلاق، محدودیت و جمع‌بندی}
شناسهٔ گوینده برای یافتن هویت واقعی استفاده نمی‌شود. نبود فرادادهٔ جمعیت‌شناختی، تحلیل عدالت بین گروه‌های جمعیتی را محدود می‌کند؛ از صدا برچسب حساس استنباط نشده است. دادهٔ انگلیسی و نسبت‌های مصنوعی انتخاب زیرمجموعه، تعمیم به فارسی یا کاربرد واقعی را تضمین نمی‌کنند.

نتیجهٔ فاز اول این است که زیرمجموعه از نظر قالب قابل پردازش است، اما تکرار دقیق، تفاوت مشارکت گویندگان و کلیپ‌های کم‌انرژی باید در طراحی تحلیل لحاظ شوند. تقسیم بدون اشتراک گوینده آماده شده و بررسی کیفیت به فهرست قابل پیگیری تبدیل شده است. کیفیت برچسب و شنیداری همهٔ نمونه‌ها تأیید نشده و هنوز ادعایی دربارهٔ توان پیش‌بینی مدل نداریم.

\newpage
\begin{thebibliography}{9}
\begin{LTR}
\bibitem{tutorial} TensorFlow. Simple audio recognition: Recognizing keywords. Accessed 2026-11-07.\\\url{https://www.tensorflow.org/tutorials/audio/simple_audio}
\bibitem{warden} Warden, P. (2018). Speech Commands: A Dataset for Limited-Vocabulary Speech Recognition. arXiv:1804.03209.\\\url{https://arxiv.org/abs/1804.03209}
\bibitem{tfds} TensorFlow Datasets. Speech Commands dataset catalog. Accessed 2026-11-07.\\\url{https://www.tensorflow.org/datasets/catalog/speech_commands}
\end{LTR}
\end{thebibliography}
\end{document}
